\cotefile{77}
\scanpage{82}
que
\[
A_{\Delta}[T]/(T^2+bT+c)
\]
est une extension quadratique étale de $A_{\Delta}$, et sa restriction à $A_{2\Delta}$ est \unsure{isom :|isom. à} $A_{2\Delta}[U]/(U^2-\Delta)$, en posant
\[
U=\frac{1}{2u}\left(T+\frac b2\right).
\]

\deleted{\unsure{Dém}}

\uline{Lemme} \deleted{\unsure{Si} \unsure{vérifie} les conditions \textcircled{a} et \textcircled{b},} \inserted{Supposons $A_{\Delta}$ int. clos dans \deleted{\unsure{\textcircled{d}}} $(A_{\Delta})_{2}$.}
% slanted
\margin{\unsure{dégager|déjà vu|dès que} ceci du \unsure{n\textsuperscript{o}|nb} 2 \uline{\unsure{comme corollaire|comme remarque|connexes}}}

Pour que $\exists\, B_{\Delta}$ ext. quadratique étale de $A_{\Delta}$, se restreignant en $A_{2\Delta}$ suivant $A_{2\Delta}[U]/(U^2-\Delta)$, il \unsure{f. et st|faut et su.} que l'on puisse trouver \inserted{$b,c\in A_{\Delta}$ et} $\unsure{\nu|v}\in\mathbb N$ et \uline{$u\in A_{\Delta}^{*}$} tels que
\[
b^2-4c=\Delta^{\deleted{\unsure{2\nu+1|\nu}}}u^2
\]

Or, quitte \unsure{à|:} multiplier $b,c$ par des puissances \inserted{\unsure{conv.|carr.|carrés}} de $\Delta$, \inserted{($\Delta^\alpha$ et $\Delta^{2\alpha}$ resp\textsuperscript{t})} \unsure{obt\textsuperscript{g}|donc} $b,c\in A$, \deleted{\unsure{Supposons}} \inserted{$u\in A$.}

\deleted{\unsure{alors}} Supposons que \deleted{\unsure{$u$}}
\[
\deleted{\unsure{\text{\textcircled{g}}}}\quad\boxed{\deleted{A_{\Delta}^{*}=\{u\Delta^d\mid d\in\mathbb Z,\ u\in A^*\}}}
\]

\deleted{donc on aura}
\[
\deleted{b^2-4c=\Delta^{2d+1}u^2\qquad(d\in\mathbb Z,\ u\in A^*).}
\]

\deleted{On a $d\in\mathbb N$ \deleted{\unsure{bien}} i.e. $d\geq 0$, car si $d<0$, on aurait}
\[
\deleted{b^2-4c=\Delta^{-\nu}u^2\qquad(\nu>0)}
\]
\deleted{donc}
\[
\deleted{\Delta^\nu(b^2-4c)=u^2}
\]
% archive sheet number: 81
